\section{Evaluation Methodology}\label{sec:eval}

The benchmark was designed before any implementation existed (\adr{019}), in
\doc{benchmark.md} for the index and \doc{evolve.md} \S11 for the optimizer.
A change to the design after the first harness run is recorded in
\doc{benchmark.md} with a date and a reason, and results from before the change
are labeled as such. Every dataset is synthetic; no real data is used, and none
is needed to reproduce a result. No result is reported yet
(Section~\ref{sec:results}).

\subsection{Questions}\label{sec:questions}

\textbf{Q1}: how many bytes does a keyed lookup read with the index, compared
with each baseline? \textbf{Q2}: how does lookup latency grow with corpus size?
\textbf{Q3}: how far is predicted cost from measured cost? Replay is exact by
design (\adr{005}), so the expected error is zero and any nonzero value is a
defect. \textbf{Q4}: given a planted anti-pattern, how often does the engine
recommend the fix an expert would make? \textbf{Q5}: how much does evolve reduce
cost per run at p50, p95, and max, at what maintenance cost, and with what gain
per replay and per model call? \textbf{Q6}: how many non-local rewrites does
evolve find that no local identity could certify, and do their output
certificates hold? \textbf{Q7}: what gain do promoted rules alone yield on
held-out query templates? \textbf{Q8}: does Pareto selection protect the p95
run relative to mean selection?

\subsection{Dataset}\label{sec:dataset}

The generator produces a synthetic versioned ledger: each entity has a true id
and a series of versions over time. Files are partitioned by day and written in
arrival order, so a version's location depends on when it arrived, not on its
entity, which imitates a common lake shape on purpose. Columns are a true id, a
version timestamp, a batch-stamp column (a high-cardinality value assigned per
write batch and ordered with time), a partition date, decimal amounts for totals
checks, and filler columns. The parameters are the number of files, row groups
per file, rows per row group, entities, churn, days, payload width, the
presence of the batch-stamp column, the fraction of null timestamps, and a
seed; the dataset is a pure function of the parameters and the seed. The
pre-registered corpus ladder is defined by days and files per day: small is 10
days of 1 file (10 files, 1,000,000 entities), medium is 30 days of 4 files
(120 files, 1,000,000 entities), and large is 30 days of 32 files (960 files,
10,000,000 entities). Every rung has 8 row groups per file, 100,000 rows per
row group, a payload width of 256 bytes, and corpus seed 7
(\doc{benchmark.md}, amended), so the configuration specifies 8, 96, and 768
million rows for the three rungs. The small corpus measured 2.1\,GB on disk;
from the configuration, the medium and large corpora would hold about 30\,GB
and 240\,GB. The ladder was revised before any harness run because its first
version, 10, 100, and 1,000 files over 30 days, was inconsistent with the
generator, which writes at least one file per day. Churn is the fraction of
entities re-versioned per day, 0.05. Rows per day are fixed by the layout
(files per day times 8 row groups of 100,000 rows), so from the configuration a
re-versioned entity would receive about 16, 64, and 51 versions per day on the
small, medium, and large rungs; the generator records the measured versions per
entity in its ground truth.
For evolve, the dataset adds a pipeline of several dependent steps
(\doc{evolve.md} \S11).

\paragraph{Planted anti-patterns.} A generated dataset may plant one or more
patterns from the rule catalog and records which ones it planted, as ground
truth beside the data: \code{batch\us stamp\us key} (consumer workloads join on
the batch-stamp column), \code{low\us selectivity\us single\us column} (a
declared single-column key where a composite prunes much better),
\code{declared\us but\us unused\us key} (a declared key spec that no workload
uses), and \code{scattered\us time\us column} (timestamps written out of order
across files).
The expert fix for each planted pattern is written down before the engine runs,
as part of the ground truth.

\subsection{Workloads and splits}\label{sec:workloads}

A recorded workload is a list of runs, and each run is a key set. Key-set sizes
are drawn to hit three points, p50, p95, and max. The other parameters are the
number of runs, the key skew (uniform or skewed draws from entities), the
recency bias (how strongly runs prefer recently changed entities), the mix of
\code{locate}, \code{latest}, and \code{asof} queries, and a seed. The
pre-registered workload has 200 recorded runs with key-set sizes of 50, 10,000,
and 100,000 at p50, p95, and max, split by run id into 150 search runs and 50
held-out runs that never enter the search, with workload seed 11
(\doc{benchmark.md}, amended; \doc{evolve.md} \S11). Runs are grouped by
template, and rule transfer is measured on held-out templates (\doc{evolve.md}
\S7).

\subsection{Baselines}\label{sec:baselines}

For lookups (\adr{019}): a raw scan that reads every row group and filters;
DuckDB running the same lookup as a query over the same files; Parquet-native
statistics, which prune with footer minimum and maximum values and any embedded
bloom filters but no sidecar index; and Apache DataFusion~\cite{datafusion}, a popular Rust engine that prunes with statistics, the page index, and embedded bloom filters. For evolve
(\doc{evolve.md} \S11): no index; Parquet statistics; DuckDB; DataFusion; an LLM-only
rewriter without replay, in the shape of GenRewrite~\cite{genrewrite}; and an
estimate-based advisor that uses a DuckDB or DataFusion cost estimate in place
of a what-if interface. Each baseline answers the same lookups on the same files
and must return the same rows as the engine. A baseline that returns different
rows is a failed run, not a data point (\doc{benchmark.md}).

\subsection{Metrics}\label{sec:metrics}

\emph{Bytes read} is measured: the harness wraps the file reader in a
byte-counting layer and records the compressed bytes actually fetched for each
run, alongside the planned bytes computed from row-group statistics
(\doc{benchmark.md}, amended). \emph{Latency versus corpus size} is the wall
time per lookup across the corpus ladder, at p50, p95, and max. An engine
baseline whose planning itself scans key columns, as Apache DataFusion's
pruning does, reports planning and execution latency separately, and realized
bytes for every strategy are counted by the harness's own reader, so an
engine's internal reads count the same way as the index's (\doc{benchmark.md},
amended).
\emph{Cost-estimation error} is
$|\mathit{predicted} - \mathit{measured}| / \mathit{measured}$, where
\emph{predicted} is the planned bytes and \emph{measured} is the counted bytes. \emph{Expert-fix recovery} is the fraction of planted anti-patterns for
which the top recommendation equals the recorded expert fix. Peak resident
memory per operation and observe-tool latency are recorded for the assertions of
\adr{018} (\doc{benchmark.md}). For evolve (\doc{evolve.md} \S11): cost
reduction per run at p50, p95, and max, with $m(a)$; blind recovery of planted
anti-patterns; realized versus predicted cost, zero for evolve by construction
and measured for the estimate-based baseline; gain per replay and per model
call; non-local rewrites found and whether their certificates held; transfer
gain from promoted rules on held-out templates; and p95 cost under Pareto versus
mean selection.

\paragraph{Planned versus measured bytes.} The realized-versus-predicted figure
plots the two byte counts against each other. The claim of no estimation error
means that their difference is zero for every strategy whose plan is executed as
written; any nonzero difference is reported, not hidden (\doc{benchmark.md},
amended). Section~\ref{sec:exact} lists the reads that a Replay does not count.

\subsection{Ablations}\label{sec:ablations}

Three ablations isolate the parts of evolve (\doc{evolve.md} \S11). Pareto
versus mean selection uses the same proposer budget and seed in both arms. The
attribution ablation replaces the exact per-node differences with the
proposer's self-reported attribution, holding everything else fixed
(\doc{benchmark.md}, amended). The gate ablation runs the search with and
without the gate; in the ungated arm every frontier member is reconciled after
the search, so the number of non-equivalent candidates the gate would have
rejected can be reported. That last detail is this paper's operationalization
and is not yet recorded in \doc{benchmark.md}.

\subsection{Environment and reproducibility}\label{sec:env}

Every result records the commit hash, the \code{rustc} version, the operating
system, the CPU model, the core count, memory, the storage kind (local disk, or
an object store and its region), and whether caches were warm or cold. Cold runs
drop the operating system's page cache before each lookup where the platform
allows; otherwise the result is labeled warm (\doc{benchmark.md}). A
configuration file holds every parameter, including seeds, and the harness
writes one JSON object per baseline and lookup. Charts are rendered from ledger
records by \code{bellwether report}, never edited by hand and never recomputed
for display (\doc{evolve.md} \S11), and a search can be replayed from its ledger
(Section~\ref{sec:determinism}).

\subsection{Results (pending)}\label{sec:results}

No result is reported yet. The harness has so far run only the small rung, on
one developer machine; no medium or large result exists, and a result shown
here will name its rung and machine. Tables~\ref{tab:lookup}
to~\ref{tab:ablation} and
Figures~\ref{fig:results-lookup} to~\ref{fig:results-loss} are the reporting
templates. Every table cell reads \emph{pending}. Every figure panel is a
pgfplots template over one CSV file, whose schema is fixed in
\doc{paper/data/README.md} and which \code{bellwether report -{}-paper}
writes from ledger records; a panel draws its file when the file exists and a
framed placeholder of the same size until then, so no panel is drawn from
anything but that file.

\begin{table*}[t]
  \caption{Lookup benchmark (Q1, Q2): bytes read and latency per lookup at the
    workload's p50, p95, and max key-set sizes, and peak resident memory.
    Pending: produced by \code{bellwether report} from ledger records.}
  \label{tab:lookup}
  \small
  \begin{tabular}{@{}lccccccc@{}}
    \toprule
    & \multicolumn{3}{c}{Bytes read} & \multicolumn{3}{c}{Latency} & Peak \\
    \cmidrule(lr){2-4} \cmidrule(lr){5-7}
    Method & p50 & p95 & max & p50 & p95 & max & memory \\
    \midrule
    Raw scan                & \pend & \pend & \pend & \pend & \pend & \pend & \pend \\
    Parquet statistics      & \pend & \pend & \pend & \pend & \pend & \pend & \pend \\
    DuckDB                  & \pend & \pend & \pend & \pend & \pend & \pend & \pend \\
    DataFusion              & \pend & \pend & \pend & \pend & \pend & \pend & \pend \\
    Bellwether index        & \pend & \pend & \pend & \pend & \pend & \pend & \pend \\
    \bottomrule
  \end{tabular}
\end{table*}

\begin{table*}[t]
  \caption{Evolve against baselines (Q5): bytes read per run at p50, p95, and
    max, and maintenance cost $m(a)$. Pending: produced by
    \code{bellwether report} from ledger records.}
  \label{tab:evolve}
  \small
  \begin{tabular}{@{}lccccc@{}}
    \toprule
    & \multicolumn{3}{c}{Bytes read per run} & \multicolumn{2}{c}{Maintenance $m(a)$} \\
    \cmidrule(lr){2-4} \cmidrule(lr){5-6}
    Method & p50 & p95 & max & Index bytes & Build time \\
    \midrule
    No index                          & \pend & \pend & \pend & \pend & \pend \\
    Parquet statistics                & \pend & \pend & \pend & \pend & \pend \\
    DataFusion                        & \pend & \pend & \pend & \pend & \pend \\
    DuckDB                            & \pend & \pend & \pend & \pend & \pend \\
    LLM-only rewriter (no replay)     & \pend & \pend & \pend & \pend & \pend \\
    Estimate-based advisor            & \pend & \pend & \pend & \pend & \pend \\
    Evolve                            & \pend & \pend & \pend & \pend & \pend \\
    \bottomrule
  \end{tabular}
\end{table*}

\begin{table}[t]
  \caption{Planted anti-patterns (Q4): the after shape from the rule catalog,
    and whether the top recommendation of \code{fix}, and of evolve run blind,
    equals the recorded expert fix. Pending.}
  \label{tab:recovery}
  \small
  \begin{tabular}{@{}>{\raggedright\arraybackslash}p{0.36\columnwidth}%
      >{\raggedright\arraybackslash}p{0.27\columnwidth}cc@{}}
    \toprule
    Pattern & Catalog after shape & \code{fix} & Evolve \\
    \midrule
    \code{batch\_stamp\_key} & join on the true key & \pend & \pend \\
    \code{low\_selectivity\_\allowbreak single\_column} & composite key &
      \pend & \pend \\
    \code{declared\_but\_\allowbreak unused\_key} & spec removed & \pend &
      \pend \\
    \code{scattered\_time\_\allowbreak column} & clustered by time & \pend &
      \pend \\
    \bottomrule
  \end{tabular}
\end{table}

\begin{table}[t]
  \caption{Cost-estimation error (Q3), non-local rewrites (Q6), and rulebook
    transfer (Q7). Pending.}
  \label{tab:error}
  \small
  \begin{tabular}{@{}>{\raggedright\arraybackslash}p{0.43\columnwidth}ccc@{}}
    \toprule
    Cost-estimation error on bytes & p50 & p95 & max \\
    \midrule
    Replay, lookups            & \pend & \pend & \pend \\
    Replay, evolve candidates  & \pend & \pend & \pend \\
    Estimate-based advisor     & \pend & \pend & \pend \\
    \midrule
    Rewrites and rules & \multicolumn{3}{c}{Count} \\
    \midrule
    Non-local rewrites that no local identity certifies
      & \multicolumn{3}{c}{\pend} \\
    Of these, output certificate held & \multicolumn{3}{c}{\pend} \\
    Non-equivalent candidates admitted without the gate
      & \multicolumn{3}{c}{\pend} \\
    Rules promoted after held-out validation & \multicolumn{3}{c}{\pend} \\
    \midrule
    Transfer gain from rules alone & p50 & p95 & max \\
    \midrule
    Held-out templates & \pend & \pend & \pend \\
    \bottomrule
  \end{tabular}
\end{table}

\begin{table}[t]
  \caption{Ablations (Q5, Q8): p95 bytes per run, gain per replay, and gain per
    model call. Pending.}
  \label{tab:ablation}
  \small
  \begin{tabular}{@{}lccc@{}}
    \toprule
    Variant & p95 bytes & Gain/replay & Gain/call \\
    \midrule
    Evolve (full)          & \pend & \pend & \pend \\
    Mean selection         & \pend & \pend & \pend \\
    Without attribution    & \pend & \pend & \pend \\
    Without the gate       & \pend & \pend & \pend \\
    \bottomrule
  \end{tabular}
\end{table}

\begin{figure*}[t]
  \centering
  \begin{subfigure}[t]{0.32\textwidth}
    \centering
    % Panel: the pruning funnel of one recorded run, one group of bars per
% strategy. Data: data/funnel.csv; schema in data/README.md.
\IfFileExists{data/funnel.csv}{%
  \begingroup\bwdatacatcodes
  % The rung the run was measured on, for the axis label.
  \pgfplotstableread{data/funnel.csv}\bwfunneltable
  \pgfplotstablegetelem{0}{files}\of\bwfunneltable
  \let\bwfunnelfiles\pgfplotsretval
  \begin{bwpanel}{4.1cm}
    \begin{axis}[bw, bw box={4.1cm}{0.95cm}{0.85cm}{0.12cm}{0.62cm},
        ybar=0pt, bar width=3.6pt, ymode=log, log origin=infty,
        xmin=-0.55, xmax=4.55, restrict x to domain=-0.5:4.5,
        xtick={0,...,4}, xmajorgrids=false,
        xticklabels={{raw\\scan}, {Parquet\\statistics}, {DuckDB}, {Data-\\Fusion},
          {Bell-\\wether}},
        xticklabel style={align=center},
        ylabel={row groups (\bwfunnelfiles\ files)},
        legend columns=-1, bw bar legend,
        every axis plot/.append style={draw=white, line width=0.3pt}]
      \addplot[fill=bwStageA]
        table[x expr=\bwstrategypos{\thisrow{strategy}}, y=rg_total] {data/funnel.csv};
      \addlegendentry{corpus}
      \addplot[fill=bwStageB]
        table[x expr=\bwstrategypos{\thisrow{strategy}}, y=rg_stats] {data/funnel.csv};
      \addlegendentry{after statistics}
      \addplot[fill=bwStageC]
        table[x expr=\bwstrategypos{\thisrow{strategy}}, y=rg_index] {data/funnel.csv};
      \addlegendentry{after index}
    \end{axis}
  \end{bwpanel}%
  \endgroup
}{%
  \pendingfigure[\scriptsize]{4.1cm}{the pruning funnel of one recorded run,
    as row groups in the corpus, after statistics, and after the index, per
    strategy}{funnel}%
}%

    \caption{Pruning funnel of one run (Q1).}
    \label{fig:funnel}
  \end{subfigure}\hfill
  \begin{subfigure}[t]{0.32\textwidth}
    \centering
    % Panel: lookup latency versus corpus size, one line per strategy at p50 with
% a band to p95. Data: data/latency.csv; schema in data/README.md.
\IfFileExists{data/latency.csv}{%
  \begingroup\bwdatacatcodes
  \begin{bwpanel}{4.1cm}
    \begin{axis}[bw, bw box={4.1cm}{0.95cm}{0.85cm}{0.12cm}{0.62cm},
        xmode=log, ymode=log, xmin=7, xmax=1400,
        xtick={10,120,960}, xticklabels={10,120,960},
        xlabel={files in corpus}, ylabel={latency per lookup (ms)},
        legend columns=3]
      \bwlatencyseries{index}{Bellwether}
      \bwlatencyseries{stats}{Parquet stats}
      \bwlatencyseries{datafusion}{DataFusion}
      \bwlatencyseries{duckdb}{DuckDB}
      \bwlatencyseries{raw}{raw scan}
    \end{axis}
  \end{bwpanel}%
  \endgroup
}{%
  \pendingfigure[\scriptsize]{4.1cm}{lookup latency versus corpus size per
    strategy, at p50 with a band to p95}{latency}%
}%

    \caption{Latency versus corpus size (Q2).}
    \label{fig:latency}
  \end{subfigure}\hfill
  \begin{subfigure}[t]{0.32\textwidth}
    \centering
    % Panel: measured versus predicted bytes per run, one marker series per
% strategy, with the line y = x. Data: data/realized.csv; schema in
% data/README.md.
\IfFileExists{data/realized.csv}{%
  \begingroup\bwdatacatcodes
  \bwloglimits{data/realized.csv}{predicted_bytes}{realized_bytes}%
  \begin{bwpanel}{4.1cm}
    \begin{axis}[bw, bw box={4.1cm}{0.95cm}{0.85cm}{0.12cm}{0.62cm},
        xmode=log, ymode=log,
        xmin=\bwlo, xmax=\bwhi, ymin=\bwlo, ymax=\bwhi,
        xlabel={predicted bytes}, ylabel={measured bytes},
        legend columns=-1]
      \addplot[bwAxis, line width=0.5pt, mark=none, forget plot]
        coordinates {(\bwlo,\bwlo) (\bwhi,\bwhi)};
      \addplot[bwA, only marks, mark=o, bw only={strategy}{index}]
        table[x=predicted_bytes, y=realized_bytes] {data/realized.csv};
      \addlegendentry{replay, lookups}
      \addplot[bwC, only marks, mark=triangle, bw only={strategy}{evolve}]
        table[x=predicted_bytes, y=realized_bytes] {data/realized.csv};
      \addlegendentry{replay, evolve}
      \addplot[bwB, only marks, mark=square, bw only={strategy}{estimate}]
        table[x=predicted_bytes, y=realized_bytes] {data/realized.csv};
      \addlegendentry{estimate}
    \end{axis}
  \end{bwpanel}%
  \endgroup
}{%
  \pendingfigure[\scriptsize]{4.1cm}{measured versus predicted bytes per run,
    per strategy, with the diagonal}{realized}%
}%

    \caption{Realized versus predicted bytes (Q3).}
    \label{fig:realized}
  \end{subfigure}
  \caption{Lookup and exactness. (a)~Row groups in the corpus, left after
    footer statistics and embedded bloom filters, and left after the sidecar
    index, for the search run at the workload's p95 key-set size on the
    largest rung run so far, whose size the axis label gives; a strategy
    without a stage passes its count through. (b)~Wall time per lookup,
    planning included: p50, with a band and a whisker up to p95. (c)~Measured
    against predicted bytes per run, with the line $y = x$; by design, replay
    lies on it, and a point off it is a defect, not noise. Every panel in
    Figures~\ref{fig:results-lookup} to~\ref{fig:results-loss} is drawn by
    pgfplots from a CSV file that \code{bellwether report -{}-paper} writes
    from ledger records (\doc{paper/data/README.md}); a framed box stands in
    for a panel whose file does not exist yet.}
  \Description{Three panels. (a) A grouped bar chart with a logarithmic axis of
    row groups: for each of five strategies, three bars for the corpus, after
    statistics, and after the index. (b) A log-log line chart of lookup latency
    in milliseconds against the number of files, one line per strategy with a
    shaded band up to the 95th percentile. (c) A log-log scatter plot of
    measured against predicted bytes per run, one marker style per strategy,
    with the diagonal. A panel whose data file does not exist is drawn as a
    framed box that names the file.}
  \label{fig:results-lookup}
\end{figure*}

\begin{figure*}[t]
  \centering
  \begin{subfigure}[t]{0.32\textwidth}
    \centering
    % Panel: convergence of evolve. Per generation, the best candidate so far and
% the generation's median candidate, each as its p95-run bytes (line) and a
% band from its p50-run to its max-run bytes. Data: data/convergence.csv;
% schema in data/README.md.
\IfFileExists{data/convergence.csv}{%
  \begingroup\bwdatacatcodes
  \begin{bwpanel}{4.1cm}
    \begin{axis}[bw, bw box={4.1cm}{0.95cm}{0.85cm}{0.12cm}{0.62cm},
        ymode=log, xlabel={generation}, ylabel={bytes per run},
        legend columns=-1]
      \addplot[name path=bwbestlo, draw=none, mark=none, forget plot]
        table[x=generation, y=best_p50] {data/convergence.csv};
      \addplot[name path=bwbesthi, draw=none, mark=none, forget plot]
        table[x=generation, y=best_max] {data/convergence.csv};
      \addplot[bwA, fill opacity=0.2, draw=none, forget plot]
        fill between[of=bwbestlo and bwbesthi];
      \addplot[name path=bwmedlo, draw=none, mark=none, forget plot]
        table[x=generation, y=median_p50] {data/convergence.csv};
      \addplot[name path=bwmedhi, draw=none, mark=none, forget plot]
        table[x=generation, y=median_max] {data/convergence.csv};
      \addplot[bwB, fill opacity=0.13, draw=none, forget plot]
        fill between[of=bwmedlo and bwmedhi];
      \addplot[bwA, solid, mark=none]
        table[x=generation, y=best_p95] {data/convergence.csv};
      \addlegendentry{best so far}
      \addplot[bwB, dashed, mark=none, unbounded coords=jump]
        table[x=generation, y=median_p95] {data/convergence.csv};
      \addlegendentry{median of generation}
    \end{axis}
  \end{bwpanel}%
  \endgroup
}{%
  \pendingfigure[\scriptsize]{4.1cm}{the best and the median candidate's cost
    per generation, with p50, p95, and max run bands}{convergence}%
}%

    \caption{Convergence (Q5).}
    \label{fig:convergence}
  \end{subfigure}\hfill
  \begin{subfigure}[t]{0.32\textwidth}
    \centering
    % Panel: Pareto frontier evolution, projected onto p50-run and p95-run bytes:
% the frontier at four generations, lighter for earlier ones, and the
% baseline. Data: data/frontier.csv; schema in data/README.md.
\IfFileExists{data/frontier.csv}{%
  \begingroup\bwdatacatcodes
  % The generation of each snapshot, for the legend.
  \pgfplotstableread{data/frontier.csv}\bwfrontiertable
  \pgfplotstableforeachcolumnelement{snapshot}\of\bwfrontiertable\as\bwsnap{%
    \pgfplotstablegetelem{\pgfplotstablerow}{generation}\of\bwfrontiertable
    \expandafter\xdef\csname bwgen\bwsnap\endcsname{\pgfplotsretval}}%
  \begin{bwpanel}{4.1cm}
    \begin{axis}[bw, bw box={4.1cm}{0.95cm}{0.85cm}{0.12cm}{0.62cm},
        xmode=log, ymode=log,
        xlabel={p50-run bytes}, ylabel={p95-run bytes},
        legend columns=-1,
        every axis plot/.append style={mark=*, mark options={solid}}]
      \addplot[bwA, opacity=0.3, bw only={snapshot}{1}]
        table[x=p50_bytes, y=p95_bytes] {data/frontier.csv};
      \addlegendentry{$g{=}\csname bwgen1\endcsname$}
      \addplot[bwA, opacity=0.53, bw only={snapshot}{2}]
        table[x=p50_bytes, y=p95_bytes] {data/frontier.csv};
      \addlegendentry{$\csname bwgen2\endcsname$}
      \addplot[bwA, opacity=0.77, bw only={snapshot}{3}]
        table[x=p50_bytes, y=p95_bytes] {data/frontier.csv};
      \addlegendentry{$\csname bwgen3\endcsname$}
      \addplot[bwA, bw only={snapshot}{4}]
        table[x=p50_bytes, y=p95_bytes] {data/frontier.csv};
      \addlegendentry{$\csname bwgen4\endcsname$}
      \addplot[bwInk, only marks, mark=star, mark size=2.4pt, line width=0.6pt,
          bw only={snapshot}{0}]
        table[x=p50_bytes, y=p95_bytes] {data/frontier.csv};
      \addlegendentry{baseline $a_0$}
    \end{axis}
  \end{bwpanel}%
  \endgroup
}{%
  \pendingfigure[\scriptsize]{4.1cm}{the Pareto frontier at four generations in
    p50-run and p95-run bytes, with the baseline}{frontier}%
}%

    \caption{Frontier evolution (Q8).}
    \label{fig:frontier}
  \end{subfigure}\hfill
  \begin{subfigure}[t]{0.32\textwidth}
    \centering
    % Panel: rulebook transfer. Held-out gain of each promoted rule applied
% alone, with its confidence as error bars, and the gain of rules 1 to i
% applied together. Data: data/transfer.csv; schema in data/README.md.
\IfFileExists{data/transfer.csv}{%
  \begingroup\bwdatacatcodes
  \begin{bwpanel}{4.1cm}
    \begin{axis}[bw, bw box={4.1cm}{1.75cm}{0.85cm}{0.12cm}{0.62cm},
        y dir=reverse, ytick=data, ymajorgrids=false,
        yticklabels from table={data/transfer.csv}{label},
        enlarge y limits={abs=0.6},
        xlabel={held-out gain (\% of bytes)},
        legend columns=-1]
      \addplot[xbar, bar width=3.6pt, fill=bwA, draw=none,
          legend image code/.code={\fill[bwA] (0cm,-0.06cm) rectangle (0.42cm,0.06cm);},
          error bars/.cd, x dir=both, x explicit,
          error bar style={bwInk, line width=0.5pt},
          error mark options={bwInk, rotate=90, mark size=1.4pt}]
        table[y=rank, x=gain_pct,
          x error plus expr=\thisrow{hi_pct}-\thisrow{gain_pct},
          x error minus expr=\thisrow{gain_pct}-\thisrow{lo_pct}]
        {data/transfer.csv};
      \addlegendentry{rule alone}
      \addplot[bwB, mark=*, mark options={solid}]
        table[y=rank, x=cumulative_pct] {data/transfer.csv};
      \addlegendentry{rules $1..i$ together}
    \end{axis}
  \end{bwpanel}%
  \endgroup
}{%
  \pendingfigure[\scriptsize]{4.1cm}{held-out gain per promoted rule, with its
    confidence, and the cumulative gain}{transfer}%
}%

    \caption{Rulebook transfer (Q7).}
    \label{fig:transfer}
  \end{subfigure}
  \caption{Search. (a)~Per generation, the best candidate so far (least
    p95-run bytes) and the generation's median candidate; lines are p95-run
    bytes and bands span p50-run to max-run bytes. (b)~The Pareto frontier at
    generations $G/4$, $G/2$, $3G/4$, and the last, $G$ (darker is later),
    projected onto p50-run and p95-run bytes, with the baseline $a_0$; the
    frontier is over every run and objective, so a member can look dominated
    in this projection. (c)~Held-out gain of each promoted rule applied alone
    (bars; error bars span its lowest to highest gain over held-out
    templates) and of rules $1$ to $i$ applied together (line).}
  \Description{Three panels. (a) A line chart of bytes per run on a
    logarithmic axis against generation, with a solid line for the best
    candidate and a dashed line for the median candidate, each with a shaded
    band. (b) A log-log plot of p95-run bytes against p50-run bytes showing
    four connected point sets, drawn lighter for earlier generations, and a
    star for the baseline. (c) A bar chart of held-out gain in percent per
    promoted rule with error bars, and a line of cumulative gain. A panel whose
    data file does not exist is drawn as a framed box that names the file.}
  \label{fig:results-search}
\end{figure*}

\begin{figure*}[t]
  \centering
  \begin{subfigure}[t]{0.27\textwidth}
    \centering
    % Panel: the loss landscape. Exact replay cost on a grid of two numeric knobs,
% bloom false-positive rate and key-count gate, as a surface in a
% three-quarter view, with the finite-difference descent path on it.
% Data: data/landscape.csv (grid) and data/descent.csv (path); schema in
% data/README.md.
\IfFileExists{data/landscape.csv}{%
  \begingroup\bwdatacatcodes
  \begin{bwpanel}{4.1cm}
    \begin{axis}[bw, bw box={4.1cm}{0.8cm}{0.65cm}{0.55cm}{0.3cm},
        view={35}{30}, xmode=log, ymode=log,
        xlabel={bloom FPR}, ylabel={key-count gate}, zlabel={bytes},
        colormap name=bwblues,
        xlabel style={sloped like x axis}, ylabel style={sloped like y axis}]
      \addplot3[surf, shader=interp, mesh/ordering=y varies]
        table[x=fpr, y=gate, z=bytes] {data/landscape.csv};
      \IfFileExists{data/descent.csv}{%
        \addplot3[bwB, mark=*, mark size=0.9pt, line width=0.7pt]
          table[x=fpr, y=gate, z=bytes] {data/descent.csv};
        \addplot3[bwB, only marks, mark=o, mark size=2pt, line width=0.6pt,
            bw only={step}{0}]
          table[x=fpr, y=gate, z=bytes] {data/descent.csv};
      }{}
    \end{axis}
  \end{bwpanel}%
  \endgroup
}{%
  \pendingfigure[\scriptsize]{4.1cm}{exact replay cost on a grid of two
    knobs as a surface, with the descent path}{landscape}%
}%

    \caption{Loss landscape and descent.}
    \label{fig:landscape}
  \end{subfigure}\hfill
  \begin{subfigure}[t]{0.235\textwidth}
    \centering
    % Panel: the same grid as data/landscape.csv as a filled contour map, with the
% descent path. Data: data/landscape.csv and data/descent.csv; schema in
% data/README.md.
\IfFileExists{data/landscape.csv}{%
  \begingroup\bwdatacatcodes
  \begin{bwpanel}{4.1cm}
    \begin{axis}[bw, bw box={4.1cm}{0.95cm}{0.92cm}{0.38cm}{0.75cm},
        view={0}{90}, xmode=log, ymode=log,
        xlabel={bloom FPR}, ylabel={key-count gate},
        colormap name=bwblues, enlargelimits=false, axis on top, grid=none,
        bw colorbar]
      \addplot3[contour filled={number=8}, mesh/ordering=y varies]
        table[x=fpr, y=gate, z=bytes] {data/landscape.csv};
      \IfFileExists{data/descent.csv}{%
        \addplot3[bwB, mark=*, mark size=0.9pt, line width=0.7pt]
          table[x=fpr, y=gate, z=bytes] {data/descent.csv};
        \addplot3[bwB, only marks, mark=o, mark size=2pt, line width=0.6pt,
            bw only={step}{0}]
          table[x=fpr, y=gate, z=bytes] {data/descent.csv};
      }{}
    \end{axis}
  \end{bwpanel}%
  \endgroup
}{%
  \pendingfigure[\scriptsize]{4.1cm}{the same grid as a contour map, with
    the descent path}{landscape}%
}%

    \caption{The same grid, contours.}
    \label{fig:contour}
  \end{subfigure}\hfill
  \begin{subfigure}[t]{0.29\textwidth}
    \centering
    % Panel: counterfactual attribution. Delta_v in bytes for each attributed
% node v of the baseline (rows) on each search run (columns), on a diverging
% scale centered on zero. Data: data/attribution.csv; schema in
% data/README.md.
\IfFileExists{data/attribution.csv}{%
  \begingroup\bwdatacatcodes
  \bwsymlimit{data/attribution.csv}{delta_bytes}%
  \bwmaxint{data/attribution.csv}{op}%
  \begin{bwpanel}{4.1cm}
    \begin{axis}[bw, bw box={4.1cm}{1.3cm}{0.92cm}{0.1cm}{0.75cm},
        colormap name=bwdiverging, point meta min=-\bwabs, point meta max=\bwabs,
        bw colorbar,
        y dir=reverse, ytick={0,...,\bwlast},
        yticklabels from table={data/attribution.csv}{label},
        xlabel={search run}, enlargelimits=false, axis on top, grid=none,
        unbounded coords=jump]
      \addplot[matrix plot*, mesh/ordering=y varies, point meta=explicit]
        table[x=run, y=op, meta=delta_bytes] {data/attribution.csv};
    \end{axis}
  \end{bwpanel}%
  \endgroup
}{%
  \pendingfigure[\scriptsize]{4.1cm}{counterfactual attribution per node and
    run}{attribution}%
}%

    \caption{Attribution $\Delta_v$.}
    \label{fig:attribution}
  \end{subfigure}\hfill
  \begin{subfigure}[t]{0.165\textwidth}
    \centering
    % Panel: pairwise interaction. I_uv in bytes, summed over search runs, for
% the nodes of the attribution panel in the same order on both axes; the
% diagonal is undefined. Data: data/interaction.csv; schema in
% data/README.md.
\IfFileExists{data/interaction.csv}{%
  \begingroup\bwdatacatcodes
  \bwsymlimit{data/interaction.csv}{i_bytes}%
  \begin{bwpanel}{4.1cm}
    \begin{axis}[bw, bw box={4.1cm}{0.2cm}{0.92cm}{0.1cm}{0.75cm},
        colormap name=bwdiverging, point meta min=-\bwabs, point meta max=\bwabs,
        bw colorbar,
        y dir=reverse, xtick distance=1, ytick distance=1,
        xticklabels={}, yticklabels={},
        xlabel={$u$, $v$: nodes of (c)}, enlargelimits=false, axis on top,
        grid=none, unbounded coords=jump,
        axis background/.style={fill=bwGrid}]
      \addplot[matrix plot*, point meta=explicit]
        table[x=v, y=u, meta=i_bytes] {data/interaction.csv};
    \end{axis}
  \end{bwpanel}%
  \endgroup
}{%
  \pendingfigure[\scriptsize]{4.1cm}{the pairwise interaction
    matrix}{interaction}%
}%

    \caption{Interaction $I_{uv}$.}
    \label{fig:interaction}
  \end{subfigure}
  \caption{The exact loss as an instrument. (a,\,b)~Exact replay bytes, the
    scalar that descent minimizes, on a grid of bloom false-positive rate and
    key-count gate, with the finite-difference descent path. This figure is
    possible only because the loss is exact and cheap enough to evaluate at
    every grid point; an estimate-based advisor has no such surface to draw.
    (c)~Attribution $\Delta_v$ (Equation~\eqref{eq:attr}) of the baseline's
    attributed nodes on each search run; red marks a node that reads more
    than its alternative would. (d)~Interaction $I_{uv}$
    (Equation~\eqref{eq:interaction}) summed over search runs; blue marks a
    pair whose joint change saves more than its two changes do separately,
    red less, and the diagonal is undefined.}
  \Description{Four panels. (a) A three-dimensional surface of bytes over two
    knobs, bloom false-positive rate and key-count gate, both on logarithmic
    axes, viewed from an oblique angle, with a descent path drawn as a line
    with markers on the surface. (b) The same grid as a filled contour map with
    the same path. (c) A heatmap of attribution in bytes with plan nodes as
    rows and search runs as columns, on a diverging blue-to-red scale. (d) A
    square heatmap of pairwise interaction in bytes on the same scale. A panel
    whose data file does not exist is drawn as a framed box that names the
    file.}
  \label{fig:results-loss}
\end{figure*}
